\subsection{Арифмитичсекие опреации над дифференцируемыми функциями}

\begin{tcolorbox}[title=Теорема 1]
    Пусть функции $f(x)$ и $g(x)$ определены в окрестности $U_a$ и дифференцируемы в точке $a$. Тогда функции $f(x) \pm g(x)$, $f(x) \cdot g(x)$ и $\frac{f(x)}{g(x)}$ (при $g(a) \neq 0$) также дифференцируемы в точке $a$, и справедливы формулы:
    
    \begin{enumerate}
        \item[(а)] $(f(x) \pm g(x))'|_a = f'(a) \pm g'(a)$
        \item[(б)] $(f(x) \cdot g(x))'|_a = f'(a)g(a) + f(a)g'(a)$
        \item[(в)] $\left( \frac{f(x)}{g(x)} \right)' \bigg|_a = \frac{f'(a)g(a) - f(a)g'(a)}{g^2(a)}$
    \end{enumerate}
\end{tcolorbox}

\begin{tcolorbox}[title=Доказательство]
    Так как $f$ и $g$ дифференцируемы, их можно представить в виде (через дифференциал):
    \[ f(x) = f(a) + f'(a)(x-a) + \bar{o}(x-a) \]
    \[ g(x) = g(a) + g'(a)(x-a) + \bar{o}(x-a) \]
    
    \textbf{а) Сумма/разность:}
    \[ f(x) \pm g(x) = (f(a) \pm g(a)) + (f'(a) \pm g'(a))(x-a) + (\bar{o}(x-a) \pm \bar{o}(x-a)) \]
    Так как сумма $\bar{o}$ есть $\bar{o}$, то выражение имеет вид линейной функции с коэффициентом $A = f'(a) \pm g'(a)$.
    Значит, производная суммы равна сумме производных.
    
    \textbf{б) Произведение:}
    Перемножим разложения:
    \[ f(x)g(x) = \left( f(a) + f'(a)(x-a) + \bar{o}(x-a) \right) \cdot \left( g(a) + g'(a)(x-a) + \bar{o}(x-a) \right) \]
    Раскрываем скобки, группируя слагаемые по степеням $(x-a)$:
    \begin{itemize}
        \item Константа: $f(a)g(a)$
        \item Линейная часть: $f(a)g'(a)(x-a) + f'(a)g(a)(x-a)$
        \item Остаток: $f'(a)g'(a)(x-a)^2 + \dots$ (все остальные слагаемые содержат $(x-a)^2$ или $\bar{o}$, то есть являются $\bar{o}(x-a)$).
    \end{itemize}
    \[ f(x)g(x) = f(a)g(a) + \underbrace{(f'(a)g(a) + f(a)g'(a))}_{A}(x-a) + \bar{o}(x-a) \]
    Следовательно, производная равна $f'(a)g(a) + f(a)g'(a)$.
    
    \textbf{в) Частное:}
    Рассмотрим приращение дроби или предел разностного отношения.
    \[ \frac{f(x)}{g(x)} - \frac{f(a)}{g(a)} = \frac{f(x)g(a) - f(a)g(x)}{g(x)g(a)} \]
    Подставим разложения для $f(x)$ и $g(x)$:
    \[ = \frac{(f(a) + f'(a)(x-a) + \bar{o})g(a) - f(a)(g(a) + g'(a)(x-a) + \bar{o})}{g(x)g(a)} \]
    В числителе $f(a)g(a)$ сокращается:
    \[ = \frac{(f'(a)g(a) - f(a)g'(a))(x-a) + \bar{o}(x-a)}{g(x)g(a)} \]
    Разделим на $(x-a)$ и перейдем к пределу $x \to a$.
    Учитывая, что $\lim_{x\to a} g(x) = g(a)$ (из непрерывности):
    \[ \lim_{x\to a} \frac{\frac{f(x)}{g(x)} - \frac{f(a)}{g(a)}}{x-a} = \frac{f'(a)g(a) - f(a)g'(a)}{g^2(a)} \]
\end{tcolorbox}

\textbf{Следствие 1 (Линейность).}
Производная линейной комбинации равна линейной комбинации производных:
\[ (c_1 f_1 + \dots + c_n f_n)' = c_1 f_1' + \dots + c_n f_n' \]

\textbf{Следствие 2 (Производная произведения $n$ функций).}
По индукции:
\[ (f_1 f_2 \dots f_n)' = f_1' f_2 \dots f_n + f_1 f_2' \dots f_n + \dots + f_1 \dots f_n' \]
(Сумма слагаемых, где в каждом штрих стоит только над одной функцией).

\textbf{Следствие 3 (В терминах дифференциалов).}
Правила для дифференциалов аналогичны:
\begin{enumerate}
    \item[(а)] $d(f \pm g) = df \pm dg$
    \item[(б)] $d(f \cdot g) = g df + f dg$
    \item[(в)] $d\left(\frac{f}{g}\right) = \frac{g df - f dg}{g^2}$
\end{enumerate}
